%% 8-0 序章　空間のどこにも矢印がある（第14回：組 V8-1。ページ制限なし。偶数ページ）
%% ① 天気図の風と砂鉄（はてな？）② 気温と風の地図を作る（やってみよう）③ 場所で決まる数・場所で決まる矢印
%% ④ はてな？の答え（8-4・8-9・8-10 の予告）⑤ この巻の地図（4つの部）⑥ この巻の約束：積分の名前と実際にすること（決定10。型は tools/kata8-int.md）
%% ⑦ 主役の山と主役の場 ⑧ 表計算と Python（型は tools/kata8.md）⑨ 新しい記号 ⑩ 路線 ⑪ この巻の読み方
%% 主役の山 h=100-0.01x^2-0.02y^2、点 (30,20) で h=83（tools/kata8.md）。主役の場は brief8-parts 3-0 の6つ。
\usetikzlibrary{decorations.markings}
\makeatletter
\gdef\bk@partlabel{第8巻}\gdef\bk@partname{空間に広がる量}
%% ---- kata8 の型（頭文字 kviiia） ----
\pgfmathdeclarefunction{kviiiah}{2}{\pgfmathparse{100-0.01*(#1)^2-0.02*(#2)^2}}
\colorlet{kviiiaX}{bkOrange!16}\colorlet{kviiiaY}{bkGreen!16}\colorlet{kviiiaP}{bkYellow!45}
\colorlet{kviiiaC}{bkGreen!60!black}
\tikzset{
  kviiia note/.style={font=\sffamily\footnotesize, text=bkRed!75!black, fill=white, inner sep=1pt, align=left},
  kviiia leader/.style={bkRed!75!black, line width=0.6pt, shorten >=1.6mm, shorten <=0.6mm}}
\pgfplotsset{
  kviiia map/.style={bkplot, axis equal image, grid=none, clip=true, axis lines=box, axis line style={-, line width=0.6pt, bkInk},
    xtick distance=20, ytick distance=20, tick align=outside, xlabel={$x$（m）}, ylabel={$y$（m）},
    every axis x label/.style={at={(axis description cs:1,0)}, anchor=north east, yshift=-5.2mm, font=\sffamily\small, text=bkInk},
    every axis y label/.style={at={(axis description cs:0,1)}, anchor=south, yshift=1.2mm, font=\sffamily\small, text=bkInk}},
  kviiia line/.style={kviiiaC, line width=0.8pt, no markers, samples=91, domain=0:360},
  kviiia bold/.style={kviiiaC, line width=1.6pt, no markers, samples=91, domain=0:360},
  kviiia bird/.style={view={30}{35}, axis lines=none, grid=none, xmin=-115, xmax=115, ymin=-85, ymax=85, zmin=0, zmax=100,
    unit vector ratio=1 1 1, clip=false},
  kviiia surf/.style={surf, shader=flat, draw=kviiiaC, line width=0.35pt, samples=11, samples y=37, domain=0:100, y domain=0:360,
    z buffer=sort, colormap={kviiia}{color=(bkGreen!28) color=(bkYellow!50) color=(bkOrange!50)}}}
\newcommand\kviiiacontour[2][kviiia line]{\addplot[#1] ({sqrt((100-(#2))/0.01)*cos(x)},{sqrt((100-(#2))/0.02)*sin(x)});}
\newcommand\kviiiaclabel[2]{\pgfmathsetmacro\kviiia@a{sqrt((100-(#1))/0.01)*cos(#2)}\pgfmathsetmacro\kviiia@b{sqrt((100-(#1))/0.02)*sin(#2)}%
  \edef\kviiia@t{\noexpand\node[fill=white, inner sep=0.8pt, font=\noexpand\sffamily\noexpand\footnotesize, text=bkGreen!35!black]
    at (axis cs:\kviiia@a,\kviiia@b) {#1};}\kviiia@t}
\newcommand\kviiiapt[2]{\node[circle, fill=bkRed, draw=white, line width=0.6pt, inner sep=0pt, minimum size=2.6mm] (kviiiaP) at (axis cs:#1,#2) {};}
\newcommand\kviiiabird[2]{%
  \draw[bkGray!70, line width=0.5pt, fill=bkGray!8] (axis cs:-115,-85,0) -- (axis cs:115,-85,0) -- (axis cs:115,85,0) -- (axis cs:-115,85,0) -- cycle;
  \draw[-{Stealth[length=2.4mm]}, line width=1pt, bkInk] (axis cs:-115,-85,0) -- (axis cs:-55,-85,0) node[anchor=north west, font=\sffamily\footnotesize, inner sep=1.5pt] {東（$x$）};
  \draw[-{Stealth[length=2.4mm]}, line width=1pt, bkInk] (axis cs:-115,-85,0) -- (axis cs:-115,-35,0) node[anchor=south east, font=\sffamily\footnotesize, inner sep=1.5pt] {北（$y$）};
  #1
  \addplot3[kviiia surf] ({sqrt(1-x/100)*100*cos(y)},{sqrt(1-x/100)*70.71*sin(y)},{x});
  #2}
\newcommand\kviiiaring[4][bkRed, line width=1.6pt]{\addplot3[#1, no markers, samples=60, samples y=0, domain=#3:#4]
  ({sqrt((100-(#2))/0.01)*cos(x)},{sqrt((100-(#2))/0.02)*sin(x)},{#2});}
\newcommand\kviiiadrop[2]{\pgfmathsetmacro\kviiia@z{kviiiah(#1,#2)}%
  \edef\kviiia@t{\noexpand\draw[bkRed!80!black, line width=0.8pt, dash pattern=on 2pt off 1.5pt] (axis cs:#1,#2,0) -- (axis cs:#1,#2,\kviiia@z);
  \noexpand\node[circle, fill=bkRed!80!black, inner sep=0pt, minimum size=1.8mm] (kviiiaQ) at (axis cs:#1,#2,0) {};
  \noexpand\node[circle, fill=bkRed, draw=white, line width=0.6pt, inner sep=0pt, minimum size=2.6mm] (kviiiaP) at (axis cs:#1,#2,\kviiia@z) {};}\kviiia@t}
\newcommand\kviiiafx[2]{\tikz[baseline=(f.base)]{\node[draw=bkLine, line width=0.6pt, fill=white, minimum height=5.5mm, inner xsep=2mm, font=\sffamily\footnotesize, minimum width=12mm] (c) {#1};
  \node[anchor=west, draw=bkLine, line width=0.6pt, fill=white, minimum height=5.5mm, inner xsep=2mm, minimum width=64mm, text width=60mm, align=left, font=\ttfamily\footnotesize] (f) at ([xshift=1mm]c.east) {#2};}}
%% ---- 「名前と操作」の型（tools/kata8-int.md。頭文字 kviiia） ----
\colorlet{kviiiaNc}{bkOrange!85!black}\colorlet{kviiiaSc}{bkTheme}\colorlet{kviiiaSb}{bkTheme!8}
\tikzset{
  kviiia dan/.style={font=\sffamily\bfseries\footnotesize, rounded corners=1.2mm, inner xsep=1.6mm, inner ysep=0.5mm, line width=0.9pt},
  kviiia dan1/.style={kviiia dan, draw=kviiiaNc, dash pattern=on 2.2pt off 1.4pt, fill=white, text=bkInk},
  kviiia dan2/.style={kviiia dan, draw=kviiiaSc, fill=kviiiaSb, text=bkInk},
  kviiia dan3/.style={kviiia dan, draw=bkTheme, fill=bkTheme, text=bkOn},
  kviiia name/.style={draw=kviiiaNc, line width=0.9pt, dash pattern=on 2.2pt off 1.4pt, rounded corners=1.2mm, inner xsep=1.2mm, inner ysep=1.2mm, fill=white},
  kviiia seki/.style={draw=kviiiaSc, line width=0.9pt, rounded corners=1.2mm, inner xsep=1.2mm, inner ysep=1.2mm, fill=kviiiaSb},
  kviiia tag/.style={font=\sffamily\bfseries\footnotesize, text=bkInk, inner sep=0.4mm}}
\newcommand\kviiiadan[1]{\tikz[baseline=(d.base)]\node[kviiia dan#1] (d) {\vphantom{あ}\ifcase#1\or ①\,名前の段\or ②\,直す段\or ③\,楽をする段\fi};}
\newcommand\kviiiadanhead[2]{\par\addvspace{8pt}\Needspace{7\baselineskip}\noindent
  \tikz[baseline=(d.base)]\node[kviiia dan#1, font=\sffamily\bfseries\small, inner ysep=0.8mm] (d) {\vphantom{あ}\ifcase#1\or ①\,名前の段\or ②\,直す段\or ③\,楽をする段\fi};%
  \hspace{0.7\zw}{\sffamily\bfseries\color{\bkTT}#2}\par\nobreak\vspace{3pt}\noindent\ignorespaces}
\newcommand\kviiiaNt[1]{\tikz[baseline=(n.base)]{\node[kviiia name] (n) {$\displaystyle #1$};\node[kviiia tag, anchor=north] at ([yshift=-0.8mm]n.south) {名前};}}
\newcommand\kviiiaSt[1]{\tikz[baseline=(n.base)]{\node[kviiia seki] (n) {$\displaystyle #1$};\node[kviiia tag, anchor=north] at ([yshift=-0.8mm]n.south) {ふつうの定積分};}}
%% ---- この章だけで使う部品（頭文字 kviiia） ----
\tikzset{
  kviiia lab/.style={font=\sffamily\footnotesize, inner sep=1pt},
  kviiia arr/.style={-{Stealth[length=1.9mm, width=1.7mm]}, line width=0.9pt},
  kviiia mid/.style={postaction={decorate}, decoration={markings, mark=at position 0.5 with {\arrow{Stealth[length=2mm, width=1.8mm]}}}}}
% 小さな矢印の場  \kviiiafield{(中心)}{vx の式}{vy の式}{色}  （\x, \y を使う。-2〜2 の格子、1目もり 6 mm、矢印は 1 につき 1.8 mm）
\newcommand\kviiiafield[4]{\begin{scope}[shift={#1}]
  \draw[bkLine, line width=0.5pt, fill=white] (-15,-15) rectangle (15,15);
  \foreach \x in {-2,...,2} {\foreach \y in {-2,...,2} {
    \pgfmathsetmacro\kviiiaU{(#2)*1.8}\pgfmathsetmacro\kviiiaW{(#3)*1.8}
    \pgfmathtruncatemacro\kviiiaOK{(abs(\kviiiaU)+abs(\kviiiaW))>0.05}
    \ifnum\kviiiaOK=1 \draw[kviiia arr, #4] ({6*\x-\kviiiaU/2},{6*\y-\kviiiaW/2}) -- ++(\kviiiaU,\kviiiaW);\else \fill[#4] (6*\x,6*\y) circle (0.5);\fi}}
\end{scope}}
% 路線の1行  \kviiiaroute{y}{路線番号}{前の駅}{この巻の駅}{この先}
\newcommand\kviiiaroute[5]{%
  \draw[bkRoute#2, line width=3pt, line cap=round] (10,#1) -- (132,#1);
  \node[inner sep=0pt] at (4,#1) {\bkrouteicon[7mm]{#2}};
  \fill[white] (24,#1) circle (2.2);\draw[bkRoute#2, line width=1.1pt] (24,#1) circle (2.2);
  \node[anchor=north, font=\sffamily\footnotesize, text=bkGray, inner sep=0pt] at (24,#1-3.8) {#3};
  \fill[white, rounded corners=1.5mm] (38,#1-4.4) rectangle (110,#1+4.4);
  \draw[bkRoute#2, line width=1.6pt, rounded corners=1.5mm] (38,#1-4.4) rectangle (110,#1+4.4);
  \node[font=\sffamily\footnotesize, text=bkInk] at (74,#1) {#4};
  \fill[bkRoute#2] (124,#1-3) rectangle (128,#1+3);
  \node[anchor=north, font=\sffamily\footnotesize, text=bkGray, inner sep=0pt] at (126,#1-3.8) {#5};}
% 部の地図の1まい  \kviiiapart{上のはしの y}{部の名前}{章}{絵}{中身}   はば132mm・高さ34mm
\newcommand\kviiiapart[5]{%
  \draw[bkTheme!80!black, line width=1pt, fill=white, rounded corners=2.5mm] (0,#1) rectangle (132,#1-34);
  \node[anchor=north west, fill=bkTheme, text=bkOn, rounded corners=1.2mm, font=\sffamily\bfseries\small,
    inner xsep=2mm, inner ysep=1mm] at (3.5,#1-3.5) {#2};
  \node[anchor=north east, font=\sffamily\footnotesize, text=bkGray, inner sep=0pt] at (128,#1-4.6) {#3};
  \begin{scope}[shift={(23,#1-21.5)}]#4\end{scope}
  \node[anchor=north west, text width=84mm, align=left, font=\footnotesize, inner sep=0pt, execute at begin node={\setlength{\baselineskip}{1.5em}}] at (44,#1-10.5) {#5};}
% 箇条書きの1項目  \kviiiaitem{文}
\newcommand\kviiiaitem[1]{\par\noindent\hangindent=1\zw\hangafter=1 ・#1\par}
% 会話の1行  \kviiiatalk{y}{sora|mina|delta}{名前}{せりふ}
\newcommand\kviiiatalk[4]{%
  \bkChar{#2}{(7,#1-6.5)}{12mm}
  \node[rectangle callout, callout absolute pointer={(15.5,#1)}, anchor=west, draw=bkTheme!80!black, line width=0.8pt,
    fill=white, rounded corners=2.5mm, inner xsep=3mm, inner ysep=1.8mm, text width=103mm, align=left,
    font=\small] at (21,#1) {#4};
  \node[anchor=north, font=\sffamily\footnotesize, text=bkGray, inner sep=0pt] at (7,#1-7.5) {#3};}
% 読み方のカード  \kviiiacard{(x,y)}{アイコン}{色}{見出し}{中身}  はば64mm・高さ54mm
\newcommand\kviiiacard[5]{%
  \node[anchor=north west, draw=#3, line width=1pt, fill=white, rounded corners=2.5mm,
    minimum width=64mm, minimum height=54mm, inner sep=0pt] at #1 {};
  % Ver2（第6巻の読者の指摘15と同じ）：外の node の anchor=north west がアイコンの中の node（∂ の字）に引きつがれて字がずれる。中は anchor=center にもどす
  \node[anchor=north west, inner sep=0pt] (ci) at ($#1+(3.5,-3)$) {\pgfkeys{/tikz/anchor=center}\bkicon[1.5em]{#2}{#3}};
  \node[anchor=west, font=\sffamily\bfseries\small, text=#3!85!black, inner sep=0pt] at ([xshift=2mm]ci.east) {#4};
  \node[anchor=north west, text width=57mm, align=left, font=\small, inner sep=0pt] at ($#1+(3.5,-12)$) {#5};}
% 場面のカード  \kviiiascene{(左上)}{章}{見出し}{絵}{説明}  はば64mm・高さ46mm
\newcommand\kviiiascene[5]{%
  \draw[bkTheme!70!black, line width=0.9pt, fill=white, rounded corners=2.5mm] #1 rectangle ++(64,-46);
  \node[anchor=north west, fill=bkTheme, text=bkOn, rounded corners=1mm, font=\sffamily\bfseries\footnotesize, inner xsep=1.6mm, inner ysep=0.8mm] at ($#1+(3,-3)$) {#2};
  \node[anchor=north west, font=\sffamily\bfseries\small, text=\bkTT, inner sep=0pt] at ($#1+(3,-10.5)$) {#3};
  \begin{scope}[shift={($#1+(16,-29)$)}]#4\end{scope}
  \node[anchor=north west, text width=29mm, align=left, font=\footnotesize, inner sep=0pt, execute at begin node={\setlength{\baselineskip}{1.5em}}] at ($#1+(33,-17)$) {#5};}
% 天気図の低気圧・高気圧  \kviiialow{(中心)}  \kviiiahigh{(中心)}
\newcommand\kviiialow[1]{\begin{scope}[shift={#1}]
  \foreach \r in {5,10,15} {\draw[bkInk!70, line width=0.7pt] (0,0) circle (\r);}
  \foreach \a in {0,45,...,315} {\draw[kviiia arr, bkBlue!80!black] ({\a}:13) -- ++({\a+120}:5.5);}
  \node[font=\sffamily\bfseries\large, text=bkRed!80!black] at (0,0) {低};
\end{scope}}
\newcommand\kviiiahigh[1]{\begin{scope}[shift={#1}]
  \foreach \r in {5,10,15} {\draw[bkInk!70, line width=0.7pt] (0,0) circle (\r);}
  \foreach \a in {0,45,...,315} {\draw[kviiia arr, bkBlue!80!black] ({\a}:7.5) -- ++({\a-60}:5.5);}
  \node[font=\sffamily\bfseries\large, text=\bkTT] at (0,0) {高};
\end{scope}}
% 棒磁石と砂鉄  \kviiiamagnet{(中心)}
\newcommand\kviiiamagnet[1]{\begin{scope}[shift={#1}]
  \clip (-24,-19) rectangle (24,19);
  \foreach \x/\y/\a in {-26.7/-17.4/54, -25.5/-13.3/45, -26.6/-9.9/33, -26.0/-6.7/24, -25.6/-3.2/12, -26.2/1.0/-4, -26.1/4.7/-17, -25.6/8.9/-31, -26.3/12.0/-40, -25.7/15.1/-49, -23.2/-16.2/57, -22.7/-13.7/51, -21.8/-9.7/40, -22.4/-5.8/24, -23.2/-2.3/9, -22.6/0.3/-1, -22.1/5.3/-23, -22.9/8.4/-33, -22.7/12.2/-46, -22.0/14.9/-56, -18.3/-16.7/71, -18.5/-12.9/59, -18.9/-9.4/45, -18.2/-6.8/36, -18.2/-2.8/15, -18.8/0.4/-2, -18.5/4.3/-23, -18.2/7.8/-40, -18.7/11.6/-54, -18.6/14.9/-65, -15.7/-16.6/80, -14.8/-13.3/73, -14.5/-10.3/63, -14.6/-6.7/46, -14.5/-2.4/19, -15.0/1.8/-13, -14.7/5.1/-36, -15.9/8.0/-48, -15.9/11.3/-62, -15.7/16.0/-78, -12.2/-17.3/96, -11.6/-12.8/86, -10.9/-9.5/77, -12.1/-6.1/54, -11.3/-3.3/37, -11.3/0.8/-10, -11.5/4.9/-49, -11.3/8.2/-69, -11.5/11.3/-81, -11.6/14.8/-92, -8.0/-16.9/117, -7.6/-13.6/112, -8.8/-9.8/92, -8.5/-5.8/78, -8.0/7.5/-90, -7.9/12.5/-107, -7.9/15.8/-116, -3.7/-16.5/148, -4.0/-12.9/139, -3.9/-10.2/136, -4.4/-6.5/124, -4.5/8.0/-126, -4.1/11.1/-136, -4.9/15.4/-136, -1.0/-16.7/172, -1.1/-13.9/169, -0.1/-10.2/179, -1.4/-6.6/162, -0.5/8.8/-173, -0.0/11.2/-180, -0.7/16.0/-173, 2.1/-16.9/-162, 2.3/-13.0/-156, 3.5/-9.7/-139, 2.0/-6.9/-154, 3.1/7.9/141, 2.6/12.2/152, 2.8/15.9/155, 6.4/-17.5/-129, 6.3/-12.8/-120, 6.5/-9.1/-108, 5.6/-6.5/-111, 7.0/4.9/90, 6.5/8.3/106, 6.7/11.5/114, 6.7/15.6/123, 10.5/-17.0/-104, 9.2/-13.3/-101, 10.6/-9.6/-80, 10.2/-6.4/-67, 10.6/5.3/58, 9.4/7.6/79, 10.6/11.6/87, 9.6/15.2/104, 13.9/-17.5/-89, 13.6/-13.1/-77, 14.0/-10.6/-66, 14.1/-6.3/-45, 13.5/-2.5/-22, 13.2/1.6/15, 13.1/4.7/39, 13.1/8.4/61, 14.2/12.4/72, 13.0/15.6/88, 17.8/-17.5/-75, 17.1/-13.1/-64, 17.6/-9.9/-51, 17.4/-5.9/-33, 17.0/-2.4/-14, 17.3/1.3/8, 17.2/4.3/25, 17.1/8.9/48, 16.7/11.2/59, 16.4/16.0/75, 21.6/-17.4/-64, 20.4/-13.9/-57, 21.0/-9.8/-42, 21.3/-6.7/-29, 21.4/-2.7/-12, 21.3/1.5/7, 21.5/4.6/20, 21.0/8.3/36, 20.3/11.8/50, 21.5/16.1/61, 23.6/-16.7/-58, 24.1/-14.0/-49, 23.8/-9.3/-35, 25.1/-6.2/-23, 24.2/-2.6/-10, 24.9/1.4/5, 23.7/4.1/16, 25.0/8.2/30, 24.7/12.2/43, 24.6/15.8/53}
    {\draw[bkInk!80, line width=0.7pt, line cap=round] (\x,\y) -- ++(\a:1.3) (\x,\y) -- ++(\a:-1.3);}
  \fill[bkBlue!70!black] (-9,-3) rectangle (0,3);\fill[bkRed!80!black] (0,-3) rectangle (9,3);
  \node[text=white, font=\sffamily\bfseries\small] at (-4.5,0) {S};\node[text=white, font=\sffamily\bfseries\small] at (4.5,0) {N};
\end{scope}}
% 見出しのアイコン（路線：線と2つの駅）
\newcommand\bk@icon@kviiiaroute{%
  \draw[bkOn, line width=0.13*\bk@iu, line cap=round] (-0.62,0) -- (0.62,0);
  \filldraw[fill=white, draw=bkOn, line width=0.09*\bk@iu] (-0.36,0) circle (0.2) (0.36,0) circle (0.2);}
\makeatother

\bkchapter{8}{0}{空間のどこにも\\矢印がある}{\bklevel{3}}{空間のどこにも矢印がある}

\noindent テレビやスマホの天気予報では、日本の地図の上に、たくさんの矢印が並んだ画面を見ることがあります。矢印は、その場所に吹いている風の向きと強さです。地図のどの場所にも、矢印が1本ずつかいてあります。

理科の実験で、棒磁石の上に紙をしいて、砂鉄をまいたことがあるでしょう。紙を軽くたたくと、砂鉄は細い線になって並び、N極からS極へ向かう模様ができます。この模様も、紙のどの場所でも「向き」が1つ決まっています。

\begin{bkfigure}
\begin{tikzpicture}[x=1mm, y=1mm]
  \fill[bkBlue!5, rounded corners=2mm] (-19,-19) rectangle (59,19);
  \draw[bkLine, line width=0.6pt, rounded corners=2mm] (-19,-19) rectangle (59,19);
  \kviiialow{(0,0)}
  \kviiiahigh{(40,0)}
  \node[anchor=north, kviiia lab] at (20,-21.5) {天気図の風（北半球）};
  \begin{scope}[shift={(90,0)}]
    \fill[white, rounded corners=2mm] (-25,-19) rectangle (25,19);
    \draw[bkLine, line width=0.6pt, rounded corners=2mm] (-25,-19) rectangle (25,19);
    \kviiiamagnet{(0,0)}
    \node[anchor=north, kviiia lab] at (0,-21.5) {棒磁石のまわりの砂鉄};
  \end{scope}
\end{tikzpicture}
\bkfigcaption{左：低気圧と高気圧のまわりの風（青い矢印）と等圧線（黒い円）。右：棒磁石のまわりに並んだ砂鉄。どちらも、場所ごとに向きが1つ決まっている。}
\label{bk:8-0:fig-kaze}
\end{bkfigure}

風の矢印も、砂鉄の並ぶ向きも、目には見えない「空気の流れ」や「磁石の力」を、場所ごとの矢印で表したものです。

\Needspace{7\baselineskip}
\begin{hatena}
天気図の風の矢印と、磁石のまわりの砂鉄の模様。どちらも「場所ごとに矢印が1本」あります。\par\vspace{3pt}
この矢印の地図から、何が読みとれるかな？\enspace 低気圧の中心や、高気圧の中心の場所は、矢印だけからわかるかな？
\end{hatena}

この巻の主役は、このような「\textbf{場所ごとに決まる量}」です。場所ごとに数が1つ決まるもの（気温・高さ）と、場所ごとに矢印が1本決まるもの（風・磁石の力）があります。序章では、まず自分で「場所ごとの量」の表を作ってみます。そのあとで、この巻の地図と約束を見ていきます。大学の内容の巻ですが、1回で分からなくても大丈夫です。図と数の表で手を動かせば、かならずわかるようになります。

\begin{yattemiyou}[気温の地図と、風の地図を作ろう]
\textbf{用意するもの}\quad 温度計（1つでよい）、メジャー、細い棒（わりばしなど）の先に \qty{20}{cm} ほどのすずらんテープをつけたもの、扇風機かエアコン
\begin{enumerate}[label={\bkcircled[bkGreen]{\arabic*}}, leftmargin=2.4\zw, itemsep=1pt, topsep=3pt]
\item 部屋の床を、たて・横 \qty{1}{m} ごとのマスに分けて、方眼の地図をノートにかく。窓のある側を地図の上にする。
\item マスのまん中、床から \qty{1}{m} の高さで気温をはかり、地図のマスに数を書く（温度計は、同じ場所に2分ほど置いてから読む）。
\item 扇風機（エアコン）をつける。同じマスで、すずらんテープの棒を床から \qty{1}{m} の高さにかざし、テープがなびく向きを矢印でかく。強くなびくほど、長い矢印にする。
\item 2つの地図をくらべよう。マスの中に書いたものは、どうちがうかな？
\end{enumerate}
\end{yattemiyou}

\noindent 次の図は、冬の朝に、窓のある部屋で、エアコン（暖房）をつけてはかった例です。左の地図のマスには\textbf{数が1つ}、右の地図のマスには\textbf{矢印が1本}入っています。

\begin{bkfigure}
\begin{tikzpicture}[x=1mm, y=1mm]
  % 気温の地図（数が1つ）
  \foreach \kviiiaR [count=\i from 0] in {{16.8,17.2,17.0,16.6},{18.1,18.6,18.9,18.2},{19.0,19.8,20.6,19.4}} {
    \foreach \v [count=\j from 0] in \kviiiaR {
      \pgfmathsetmacro\kviiiaC{(\v-16)*14}
      \fill[bkOrange!\kviiiaC] (12*\j,-10*\i) rectangle ++(12,-10);
      \node[font=\sffamily\small] at (12*\j+6,-10*\i-5) {\v};}}
  \draw[bkInk!60, line width=0.5pt, xstep=12, ystep=10] (0,-30) grid (48,0);
  \draw[bkBlue!70!black, line width=2pt] (8,0.8) -- (40,0.8);
  \node[anchor=south, kviiia lab, text=bkBlue!70!black] at (24,2) {窓};
  \node[anchor=north, kviiia lab] at (24,-32) {気温（℃）：マスに数が1つ};
  % 風の地図（矢印が1本）
  \begin{scope}[shift={(72,0)}]
    \draw[bkInk!60, line width=0.5pt, xstep=12, ystep=10] (0,-30) grid (48,0);
    \draw[bkBlue!70!black, line width=2pt] (8,0.8) -- (40,0.8);
    \node[anchor=south, kviiia lab, text=bkBlue!70!black] at (24,2) {窓};
    \fill[bkGray!60] (36,-31.2) rectangle (47,-29.6);
    \node[anchor=north, kviiia lab] at (41.5,-32) {エアコン};
    \foreach \i/\j/\kviiiaU/\kviiiaW in {0/0/-1/1.4, 0/1/0/2, 0/2/0/2.4, 0/3/0.6/1.8, 1/0/-1.8/0.8, 1/1/-1.4/2.2, 1/2/-0.6/3.2, 1/3/0/3.0, 2/0/-2.4/0.4, 2/1/-3.2/0.6, 2/2/-3.6/1.6, 2/3/-1.2/4.2} {
      \draw[kviiia arr, bkBlue!80!black] ({12*\j+6-1.1*\kviiiaU},{-10*\i-5-1.1*\kviiiaW}) -- ++({2.2*\kviiiaU},{2.2*\kviiiaW});}
    \node[anchor=north, kviiia lab] at (12,-32) {風：マスに矢印が1本};
  \end{scope}
\end{tikzpicture}
\bkfigcaption{はかった例。左は気温（オレンジがこいほど暖かい）、右はすずらんテープがなびいた向き（矢印が長いほど強い）。エアコンは右下のかべの上のほうにある。}
\label{bk:8-0:fig-heya}
\end{bkfigure}

左の地図からは、窓に近い所ほど気温が低いことが読みとれます。右の地図からは、エアコンの風が上（窓のほう）へ向かい、左へ回っていくようすが読みとれます。マスの大きさを小さくしていくと、部屋の中の\textbf{すべての場所}に、数や矢印が決まっていると考えられます。

\begin{omoidasou}
\bkrecall{3-2}{場所は、2つの数の組（座標）$(x,\ y)$ で表せる。}
\bkrecall{3-5}{グラフの傾きは「変化の速さ」。地図の等高線の間がせまい所ほど、坂は急。「東へ進むときと北へ進むときの傾きは、8-2で学びます」。}
\bkrecall{4-10}{向きと大きさをもつ量がベクトル（矢印 $\vec F$）。大きさだけで決まる量がスカラー。}
\bkrecall{4-13}{空間の点は3つの数 $(x,\ y,\ z)$ で表せる。3次元のベクトルの成分も3つ。}
\bkrecall{5-10}{曲がった形の面積は、細い長方形に分けてたし、細かくした先の値（定積分）で求める。}
\bkrecall{7-7}{座標の軸を回しても、矢印そのものは変わらない。成分の値だけが変わる。}
\end{omoidasou}

%% ---------------- ① 場所で決まる数・場所で決まる矢印 ----------------
\Needspace{12\baselineskip}\bksection{doguu}{場所ごとに決まる量}

\bksubsection{1}{場所で決まる数と、場所で決まる矢印}

\label{bk:8-0:ba}やってみようの2つの地図のように、身のまわりには、「場所を決めると、何かが1つ決まる」量がたくさんあります。決まるものが数1つか、矢印1本かで、2つの仲間に分けられます。

\begin{bkcenter}
\begin{bktable}{colspec={Q[l,m,wd=58mm] Q[l,m,wd=62mm]}, rowsep=1.6pt, cells={font=\small}, row{1}={font=\sffamily\bfseries\small, bg=bkTheme!10}}
場所を決めると、\textbf{数が1つ}決まる & 場所を決めると、\textbf{矢印が1本}決まる \\
山の高さ（標高）、気温、気圧 & 風の速さと向き（風速） \\
部屋の中の温度、水の深さ & 川・管の水の流れの速さと向き \\
電位（電気の高さ）、明るさ & 電場（電気の力）、磁場（磁石の力） \\
ものの密度（\qty{1}{m^3} あたりの質量） & 熱の流れ、重力 \\
\end{bktable}
\bktabcaption{場所で決まる量の、2つの仲間}
\end{bkcenter}

場所ごとに数が1つ決まっている量を\textbf{スカラー場}（スカラーば）、場所ごとに矢印が1本決まっている量を\textbf{ベクトル場}（ベクトルば）といいます。\bkidx{べくとるば}{ベクトル場}\bkidx{すからーば}{スカラー場}「場」は、「場所ごとに決まっている量の全体」という意味のことばです。スカラーとベクトルは、4-10 で学んだことばです（大きさだけの量と、向きと大きさをもつ量）。

スカラー場は、場所を入れると数が1つ出てくる「入り口が2つ（または3つ）の関数」です。地図の上なら、東へ $x$〔m〕、北へ $y$〔m〕の場所を入れると、高さ $h$〔m〕が出てきます。8-1 で、この関数を地図の等高線と表で見ていきます。ベクトル場は、場所を入れると矢印が出てくる関数です。8-5 で、矢印の図にかいて見ていきます。

\begin{bkfigure}
\begin{tikzpicture}[x=1mm, y=1mm]
  \node[draw=\bkTT, line width=1pt, fill=bkTheme!8, rounded corners=2mm, minimum width=30mm, minimum height=14mm, font=\sffamily\small, align=center] (m1) at (0,0) {高さの\\関数};
  \node[anchor=east, font=\small] (i1) at ([xshift=-12mm, yshift=3.5mm]m1.west) {東へ $x$};
  \node[anchor=east, font=\small] (i2) at ([xshift=-12mm, yshift=-3.5mm]m1.west) {北へ $y$};
  \draw[bk arrow] ([xshift=1.5mm]i1.east) -- ([yshift=3.5mm, xshift=-1.5mm]m1.west);
  \draw[bk arrow] ([xshift=1.5mm]i2.east) -- ([yshift=-3.5mm, xshift=-1.5mm]m1.west);
  \node[anchor=west, font=\small] (o1) at ([xshift=12mm]m1.east) {高さ $h$（数1つ）};
  \draw[bk arrow] ([xshift=1.5mm]m1.east) -- ([xshift=-1.5mm]o1.west);
  \begin{scope}[shift={(0,-24)}]
  \node[draw=\bkTT, line width=1pt, fill=bkTheme!8, rounded corners=2mm, minimum width=30mm, minimum height=14mm, font=\sffamily\small, align=center] (m2) at (0,0) {風の\\関数};
  \node[anchor=east, font=\small] (j1) at ([xshift=-12mm, yshift=3.5mm]m2.west) {東へ $x$};
  \node[anchor=east, font=\small] (j2) at ([xshift=-12mm, yshift=-3.5mm]m2.west) {北へ $y$};
  \draw[bk arrow] ([xshift=1.5mm]j1.east) -- ([yshift=3.5mm, xshift=-1.5mm]m2.west);
  \draw[bk arrow] ([xshift=1.5mm]j2.east) -- ([yshift=-3.5mm, xshift=-1.5mm]m2.west);
  \node[anchor=west, font=\small] (o2) at ([xshift=12mm]m2.east) {風の矢印 $\vec v$（矢印1本）};
  \draw[bk arrow] ([xshift=1.5mm]m2.east) -- ([xshift=-1.5mm]o2.west);
  \end{scope}
\end{tikzpicture}
\bkfigcaption{場所を入れると、数が1つ出てくる関数（上）と、矢印が1本出てくる関数（下）。3-3 の「入れると出てくる箱」の、入り口が2つになったもの。}
\end{bkfigure}

\bksubsection{2}{矢印が、空間のすべての点に}

\label{bk:8-0:kuukan}第4巻では、1本の矢印（ベクトル）を、たし算したり、内積・外積を計算したりしました。第7巻では、矢印を行列でのばしたり、回したりしました。どちらも、考えていた矢印は1本か2本でした。

この巻では、\textbf{空間のすべての点に、矢印が1本ずつ}ある場面を考えます。部屋の空気なら、床の近くにも、天井の近くにも、部屋のまん中にも、風の矢印があります。点の場所は3つの数 $(x,\ y,\ z)$ で表せるので（4-13）、矢印は3つの数を入れると決まる関数です。矢印の成分も3つあるので、ベクトル場は「3つの数を入れると、3つの数が出てくる関数」です。
\[ \vec F(x,\ y,\ z)=\bigl(F_x(x,\ y,\ z),\ F_y(x,\ y,\ z),\ F_z(x,\ y,\ z)\bigr) \]
たとえば、$F_x(x,\ y,\ z)$ は、「点 $(x,\ y,\ z)$ にある矢印の、$x$ の向きの成分」です。いまは、この式の形を覚えなくてもだいじょうぶです。8-5 で、まず平面（成分が2つ）でていねいに見てから、3次元に広げます。

%% ---------------- ② はてな？の答え ----------------
\bksubsection{3}{矢印の地図から読みとれること}

\label{bk:8-0:kotae}\noindent\textbf{\bkhatenaref の答え}\par
図\ref{bk:8-0:fig-kaze}の左の天気図では、低気圧のまわりの風が、左回りにうずを巻きながら中心へ吹きこんでいます。高気圧からは、右回りに外へ吹き出しています。右の砂鉄は、N極から出て、まわりこんでS極へ入る線に並んでいます。矢印の地図から、とくに次の3つのことが読みとれます。

\begin{bkfigure}
\begin{tikzpicture}[x=1mm, y=1mm]
  \kviiiafield{(0,0)}{\x}{\y}{bkOrange!85!black}
  \node[anchor=north, kviiia lab, align=center] at (0,-17) {\textbf{わき出す所}\\矢印が外へ広がる\\（8-9 発散 div）};
  \kviiiafield{(46,0)}{-\y}{\x}{bkPurple!85!black}
  \node[anchor=north, kviiia lab, align=center] at (46,-17) {\textbf{うずを巻く所}\\矢印がぐるりと回る\\（8-10 回転 rot）};
  \begin{scope}[shift={(92,0)}]
    \draw[bkLine, line width=0.5pt, fill=white] (-15,-15) rectangle (15,15);
    \begin{scope}
    \clip (-15,-15) rectangle (15,15);
    \foreach \k in {1,...,5} {\draw[kviiiaC, line width=0.8pt] ({5-1.4*\k},0) circle ({2.8*\k});}
    \fill[kviiiaC] (5,0) circle (0.6);
    \end{scope}
    \draw[kviiia arr, bkRed!80!black, line width=1pt] (-6.6,0) -- ++(-2.2,0);
    \draw[kviiia arr, bkRed!80!black, line width=1pt] (9.2,0) -- ++(5,0);
    \draw[kviiia arr, bkRed!80!black, line width=1pt] (5,-5.6) -- ++(0,-3.4);
  \end{scope}
  \node[anchor=north, kviiia lab, align=center] at (92,-17) {\textbf{等高線がつまった所}\\急な向きに長い矢印\\（8-4 勾配）};
\end{tikzpicture}
\bkfigcaption{矢印の地図から読みとれる3つのこと。「等高線がつまった所」の絵は、頂上（緑の点）が右によった山の等高線。赤い矢印はボールの転がる向き（勾配 grad の反対の向き）で、線の間がせまい右ほど坂が急で、矢印が長い。}
\label{bk:8-0:fig-yomu}
\end{bkfigure}

\begin{enumerate}[label={\bkcircled[\bkTT]{\arabic*}}, leftmargin=2.4\zw, itemsep=2pt, topsep=2pt]
\item \textbf{矢印がわき出す所・吸いこまれる所}：高気圧の中心（空気が上からおりてきて、まわりへ広がる）や、プラスの電気をもったもの（電気の力の線がわき出す）、水をそそぐ口のまわり。8-9 で、「小さな箱から出ていく量」として数で表します（発散）。
\item \textbf{矢印がうずを巻く所}：低気圧や台風の中心。8-10 で、「小さな風車が回るかどうか」として数で表します（回転）。
\item \textbf{いちばん急な向き}：天気図の等圧線（気圧の等高線）のつまった所では、気圧の差が大きく、風が強くなります。8-4 で、「坂がいちばん急な向き」の矢印として表します（勾配）。
\end{enumerate}

ただし、電荷や水をそそぐ口のように、矢印がわき出す・吸いこまれる「点」を見つけるには、矢印をよく見るだけでなく、まわりを囲んで出入りを数える方法が要ります。それが、8-9 と 8-11 の道具です。砂鉄の線は、N極から出てS極へ入るように見えますが、じつは磁石の中を通ってひとまわりし、輪になっています。磁石の力の矢印には、わき出しも吸いこみもありません（N極だけの磁石はない。8-9・8-14）。

\Needspace{9\baselineskip}
\begin{seikaku}
天気図の風は、等圧線と直角に（気圧の高い所から低い所へまっすぐ）は吹きません。地球が回っているため、北半球では、動く空気は進む向きの右へそらされます（4-14 のコラム、コリオリの力）。そのため、上空の風は等圧線にほぼ平行に吹き、地上の風は、地面とのまさつで少し低気圧のほうへ曲がって、等圧線をななめに横切ります。図\ref{bk:8-0:fig-kaze}の矢印は、この地上の風のようすをかいたものです。8-4 では、まず「気圧の傾きが空気を押す力」を、等圧線と直角な矢印として求めます。
\end{seikaku}

\begin{column}{天気予報のコンピュータは、何を計算している？}
天気予報のコンピュータは、地球をつつむ空気を、たて・横・高さに細かいマス（格子）に分けています。マスの1つ1つに、気温・気圧・水蒸気の量という「数」と、風という「矢印」を持たせます。つまり、空気を、この章で見た「スカラー場」と「ベクトル場」として、数の表にしているのです。

そして、「となりのマスとの気圧の差で、風がどう変わるか」「風に運ばれて、気温や水蒸気がどう動くか」を表す式を使って、少し先の時刻のマスの値を計算します。これを何千回もくり返して、明日やあさっての天気図を作ります。この計算に使う式は、場所と時間の両方で変わる量の式（偏微分方程式）です。8-13 では、同じ考え方で、金属の棒に熱が広がるようすを表計算で計算します。
\end{column}

\bksubsection{4}{ベクトル解析が働く場面}

\label{bk:8-0:bamen4}場所ごとの数や矢印を計算する道具は、身のまわりの製品を作るときに、毎日のように使われています。この巻で、次のような場面に何回も出会います。

\begin{bkfigure}
\begin{tikzpicture}[x=1mm, y=1mm]
  \kviiiascene{(0,0)}{8-4・8-13}{スマホの熱をにがす}{%
    \fill[bkGray!35, rounded corners=0.8mm] (-12,-8) rectangle (12,-4);
    \filldraw[fill=bkOrange!60, draw=bkOrange!80!black, line width=0.6pt] (-4,-4) rectangle (4,0);
    \foreach \a in {30,60,90,120,150} {\draw[kviiia arr, bkRed!80!black] (0,1) ++(\a:2) -- ++(\a:6);}}%
    {熱い部品から、熱がどの向きにどれだけ流れるか（温度の勾配）を計算する。}
  \kviiiascene{(68,0)}{8-5・8-14}{車のまわりの風}{%
    \foreach \y/\b in {-1/2.6, 3/2.2, 7/1.2} {\draw[bkBlue!80!black, line width=0.7pt] (-14,\y) .. controls (-4,{\y+\b}) and (4,{\y+\b}) .. (14,\y);}
    \filldraw[fill=bkRed!45, draw=bkRed!80!black, line width=0.7pt, rounded corners=1.2mm] (-9,-9) -- (-9,-5) -- (-4,-5) -- (-1,-1.5) -- (5,-1.5) -- (8,-5) -- (10,-6) -- (10,-9) -- cycle;
    \fill[bkInk] (-5,-9.3) circle (1.4) (6,-9.3) circle (1.4);}%
    {空気の流れの矢印を計算して、風の力が小さい形をさがす。}
  \kviiiascene{(0,-50)}{8-10・8-11}{モーターと発電機}{%
    \foreach \r in {3.5,6.5,9.5} {\draw[bkPurple!80!black, line width=0.7pt] (0,0) circle (\r);}
    \draw[kviiia arr, bkPurple!80!black] (9.5,0) arc (0:40:9.5);
    \filldraw[fill=bkOrange!40, draw=bkOrange!80!black, line width=0.6pt] (0,0) circle (1.7);
    \fill[bkInk] (0,0) circle (0.5);}%
    {電流のまわりにできる磁場のうず（回転）で、回る力や電気を生み出す。}
  \kviiiascene{(68,-50)}{8-4・8-12}{半導体の中の電気}{%
    \foreach \i/\kviiiaC in {0/70,1/55,2/42,3/30,4/20,5/12} {\fill[bkBlue!\kviiiaC] ({-12+4*\i},-6) rectangle ++(4,12);}
    \draw[bkInk, line width=0.6pt] (-12,-6) rectangle (12,6);
    \draw[kviiia arr, bkRed!80!black, line width=1pt] (-8,0) -- (8,0);
    \node[kviiia lab, anchor=south] at (-10,6.6) {高};\node[kviiia lab, anchor=south] at (10,6.6) {低};}%
    {電気の高さ（電位）の地図から、電子を動かす電場を計算する。}
\end{tikzpicture}
\bkfigcaption{この巻の道具が働く4つの場面。「モーターと発電機」の絵は、紙の面から手前へ流れる電流（まん中の点）のまわりの磁場。}
\end{bkfigure}

%% ---------------- ③ この巻の地図 ----------------
\bksection{goal}{この巻の地図}

\bksubsection{5}{4つの部}

\label{bk:8-0:chizu}第8巻は、4つの部に分かれています。どの章も大学（工学部の1〜2年）で学ぶ内容です。★★★★★の章は、「目でわかる」「手でわかる」「式でわかる」の3段に分けてあり、「ここまでで十分」の旗の所まで読めば、その章のゴールです（第6巻・第7巻と同じ）。

\begin{bkfigure}
\begin{tikzpicture}[x=1mm, y=1mm]
  \kviiiapart{0}{第1部\enspace たくさんの変数の微分}{8-1〜8-4}%
    {\foreach \a/\b in {4/2.8, 9/6.3, 14/9.8} {\draw[kviiiaC, line width=0.8pt] (0,0) ellipse ({\a} and {\b});}
     \draw[kviiia arr, bkRed!80!black, line width=1.1pt] (10.6,6.0) -- ++(-5.2,-4.2);}
    {\kviiiaitem{場所で決まる高さを、地図の等高線と表で読む（8-1）。}\kviiiaitem{1つだけ動かしたときの変わり方（偏微分 8-2）、ななめに少し動いたときの変化（全微分 8-3）。}\kviiiaitem{いちばん急な向きの矢印（勾配 grad 8-4）。}}
  \kviiiapart{-37}{第2部\enspace ベクトル場と積分}{8-5〜8-8}%
    {\foreach \y in {-6,0,6} {\foreach \x in {-12,-4,4} {\draw[kviiia arr, bkBlue!80!black] (\x,\y) -- ++(5,{0.9*sin(20*\y+30*\x)});}}
     \draw[bkOrange!85!black, line width=1.6pt] (12,-9) -- (12,9);}
    {\kviiiaitem{場所ごとの矢印（ベクトル場 8-5）。}\kviiiaitem{面や立体でたす（重積分 8-6）。道にそってたす（線積分 8-7）。}\kviiiaitem{面を通りぬける量をたす（面積分 8-8）。}}
\end{tikzpicture}
\bkfigcaption{第8巻の4つの部（その1）。第1部と第2部では、微分と積分の道具を作る。}
\end{bkfigure}

第1部で作る「1つだけ動かして変わり方を見る」道具（偏微分）は、第3部の発散・回転の計算にそのまま使います。第2部で作る「タイルに分けてたす」道具（積分）は、第3部の2つの定理で、発散・回転とつながります。

\begin{bkfigure}
\begin{tikzpicture}[x=1mm, y=1mm]
  \kviiiapart{-74}{第3部\enspace 発散と回転}{8-9〜8-12}%
    {\foreach \a in {0,60,...,300} {\draw[kviiia arr, bkOrange!85!black] (-8,0) ++(\a:2.2) -- ++(\a:4.2);}
     \foreach \a in {0,60,...,300} {\draw[kviiia arr, bkPurple!85!black] (9,0) ++(\a:5) -- ++(\a+90:4);}}
    {\kviiiaitem{わき出し（発散 div 8-9）と、うず（回転 rot 8-10）。}\kviiiaitem{中のわき出しと外へ出る量をつなぐ定理（8-11）。}\kviiiaitem{まわりの平均との差（ラプラシアン 8-12）。}}
  \kviiiapart{-111}{第4部\enspace 偏微分方程式}{8-13〜8-14}%
    {\foreach \i/\kviiiaC in {0/90,1/70,2/52,3/38,4/27,5/20} {\fill[bkOrange!\kviiiaC] ({-14+4.6*\i},-6) rectangle ++(4.6,6);}
     \draw[bkInk, line width=0.6pt] (-14,-6) rectangle (13.6,0);
     \draw[bkInk, line width=0.8pt, domain=-14:13.6, samples=40] plot (\x,{4+2*sin(30*\x)});}
    {\kviiiaitem{棒に広がる熱、弦を伝わる波の式を、表計算で解く（8-13）。}\kviiiaitem{電気と磁気の4つの式・流れの式を、「わき出し」と「うず」のことばで読む（8-14）。}}
\end{tikzpicture}
\bkfigcaption{第8巻の4つの部（その2）。第3部と第4部では、作った道具を組み合わせて、物理の式を読む。}
\end{bkfigure}

第4部では、それまでの道具を全部使って、熱や波や電気の式を読みます。8-14 を読み終えると、電気と磁気の基本の式（マクスウェル方程式）が、「どこでわき出し、どこでうずを巻くか」を述べた式として読めるようになります。

%% ---------------- ④ この巻の約束 ----------------
\bksection{hon}{この巻の約束：積分の「名前」と、実際にすること}

\bksubsection{6}{記号は「操作の名前」。計算は、ふつうの定積分で}

\label{bk:8-0:yakusoku}第5巻の定積分は、細い長方形に分けてたし、細かくした先の値でした。この巻では、たし合わせる場所が\textbf{面・立体・道・曲がった面}に広がり、新しい積分の記号が出てきます。

\begin{bkcenter}
\begin{bktable}{colspec={Q[c,m,wd=40mm] Q[l,m,wd=80mm]}, rowsep=1.2pt, cells={font=\small}, row{1}={font=\sffamily\bfseries\small, bg=bkTheme!10}}
名前（記号） & ふつうの定積分に直すと \\
$\displaystyle\iint_D f\,\dd A$\enspace $\displaystyle\iiint_V f\,\dd V$ & 場所を2つ（3つ）の数で表し、定積分を2回（3回）（8-6） \\
$\displaystyle\int_C\vec F\cdot\dd\vec r$ & 道を1つの数で表し、定積分を1回（8-7） \\
$\displaystyle\iint_S\vec F\cdot\vec n\,\dd S$ & 面を2つの数で表し、重積分に（8-8） \\
$\displaystyle\oint_C$\enspace $\displaystyle\oiint_S$ & 上と同じ。◯は「閉じている」という印（8-7・8-8・8-11） \\
\end{bktable}
\bktabcaption{積分の名前と、ふつうの定積分に直した形（くわしくは、それぞれの章で）}
\end{bkcenter}

これらの記号は、どれも「\textbf{タイルに分けて、1枚ずつ量をかけて、たし合わせ、タイルを細かくした先}」という\textbf{操作の名前}です。記号を見ただけでは、まだ計算できません。実際の値は、場所を数で表して、第5巻の\textbf{ふつうの定積分}（のくり返し）に直して求めます。$\int$ の数は場所を表す数の数（道は1つ、面は2つ、立体は3つ）、◯は円の形ではなく「1周する・すっぽりつつむ」という印です。この巻では、点線の枠を「名前」、実線の枠を「ふつうの定積分」の印にします（8-6 の板の重さ）。
\[ \kviiiaNt{\iint_D\sigma\,\dd A}=\kviiiaSt{\int_0^1\!\!\Bigl(\int_0^2(10+5x+2y)\,\dd x\Bigr)\dd y} \]
どう直すかは、面や道の形によって場合ごとにちがいます。そこで、積分の章（8-6・8-7・8-8・8-11）では、どの積分も次の3つの段で考え、「こんなときは、こうする」という\textbf{場合分けの表}と\textbf{流れ図}を置きます。巻末の「積分の地図」に、全部の場合をまとめます。\par\vspace{3pt}
\noindent\kviiiadan{1}\enspace 何を・どこで・1枚の量は何か、を記号から読む。\par\vspace{2pt}
\noindent\kviiiadan{2}\enspace 場所を数で表し、範囲を決めて、ふつうの定積分を1〜3回計算する。\par\vspace{2pt}
\noindent\kviiiadan{3}\enspace 楽な道（対称性・分ける・定理・表計算）をさがす。閉じた面や道は、答えが変わらない範囲で、計算が楽な形に取りかえてよい（8-11）。②の計算を始める前に確かめる。

%% ---------------- ⑤ 主役の山と主役の場 ----------------
\bksection{butsuri}{この巻の主役}

\bksubsection{7}{主役の山}

\label{bk:8-0:yama}この巻の前半には、主役の山が出てきます。頂上の高さは \qty{100}{m} で、頂上から東へ $x$〔m〕、北へ $y$〔m〕の所の高さ $h$〔m〕が、次の式で決まる山です。
\[ h=100-0.01x^2-0.02y^2 \]
たとえば、頂上から東へ \qty{30}{m}、北へ \qty{20}{m} の点 $(30,\ 20)$ の高さは、$x=30$、$y=20$ を入れて
\[ h=100-0.01\times30^2-0.02\times20^2=100-9-8=83 \]
なので、\qty{83}{m} です。$0.01x^2$ と $0.02y^2$ は「頂上からはなれるほど、どれだけ低くなるか」を表します。$y^2$ の前の数のほうが大きいので、北や南へ進むほうが、東や西へ進むより早く低くなります（北の坂のほうが急）。

\begin{bkfigure}
\begin{tikzpicture}
\begin{axis}[kviiia bird, width=100mm]
  \kviiiabird{}{\kviiiaring{80}{166}{418}\kviiiadrop{30}{20}
    \node[kviiia note, anchor=west] (n) at (axis cs:62,20,98) {点 $(30,\ 20)$ の上\\高さ 83 m};
    \draw[kviiia leader] (n.west) -- (kviiiaP);
    \node[kviiia note, anchor=west, text=bkRed!75!black] (m) at (axis cs:140,-45,62) {高さ 80 m の線};
    \draw[kviiia leader] (m.west) -- (axis cs:22.36,-27.39,80);}
\end{axis}
\end{tikzpicture}
\bkfigcaption{主役の山を、南東の空から見下ろした図（鳥かん図）。横の輪は 10 m ごとの高さの線。赤い輪は高さ 80 m の線。点線は、山の中をかくれて通る線。}
\label{bk:8-0:fig-yama}
\end{bkfigure}

この1つの山を、章ごとにちがう目で見ていきます。

\begin{bkcenter}
\begin{bktable}{colspec={Q[c,m,wd=11mm] Q[l,m,wd=30mm] Q[l,m,wd=86mm]}, rowsep=1.2pt, cells={font=\small}, row{1}={font=\sffamily\bfseries\small}}
章 & 見る目 & 主役の山の点 $(30,\ 20)$ では \\
8-1 & 地図と等高線 & 高さ \qty{83}{m}。等高線はだ円になる \\
8-2 & 東の坂・北の坂 & 東へ \qty{1}{m} で約 \qty{0.6}{m}、北へ約 \qty{0.8}{m} 下がる \\
8-3 & ななめに少し動く & 東と北へ \qty{1}{m} ずつで、約 \qty{1.4}{m} 下がる \\
8-4 & いちばん急な向き & 等高線に直角な向きがいちばん急（坂は1）。頂上への向きとは少しずれる \\
8-6 & 山の体積 & 山全体の土の量は、約 $\qty{1.11e6}{m^3}$ \\
8-7 & 登る仕事 & \qty{50}{kg} の人が登る仕事は、どの道でも \qty{8330}{J} \\
8-12 & まわりとの差 & どこでも「まわりの平均より高い」（山の形） \\
\end{bktable}
\bktabcaption{主役の山の、章ごとの顔（くわしくは、それぞれの章で）}
\end{bkcenter}

\bksubsection{8}{主役の6つの場}

\label{bk:8-0:ba6}この巻の後半には、平面の上の6つのベクトル場が、主役として何回も出てきます。どれも、点 $(x,\ y)$ に矢印 $\vec F=(F_x,\ F_y)$ を1本ずつ置いたものです。

\begin{bkfigure}
\begin{tikzpicture}[x=1mm, y=1mm]
  \foreach \cx/\cy/\fx/\fy/\nm/\ff in {0/0/{\x}/{\y}/（ア）広がる/{$(x,\ y)$}, 46/0/{-\x}/{-\y}/（イ）集まる/{$(-x,\ -y)$}, 92/0/{1.5}/{0}/（ウ）そろった/{$(1,\ 0)$},
      0/-44/{-\y}/{\x}/（エ）うず/{$(-y,\ x)$}, 46/-44/{\y}/{0}/（オ）ずれ/{$(y,\ 0)$}, 92/-44/{\x}/{0}/（カ）側溝/{$(x,\ 0)$}} {
    {\kviiiafield{(\cx,\cy)}{\fx}{\fy}{\bkTT}}
    \node[anchor=north, kviiia lab, align=center] at (\cx,\cy-16.5) {\textbf{\nm}\enspace\ff};}
\end{tikzpicture}
\bkfigcaption{主役の6つの場（矢印の長さは縮めてかいた）。8-5 で絵にし、8-7・8-9・8-10・8-11 で同じ場を使う。}
\label{bk:8-0:fig-ba6}
\end{bkfigure}

「うず」の場 $(-y,\ x)$ と、それによく似た $(y,\ x)$ は、物理シリーズ第1巻（力学）の仕事の章（力-4）にも出てきます。どちらも同じように回っているように見えますが、8-7 では、遠回りすると仕事が変わるかどうかで、2つがまったくちがう場であることを見ます。

\bksubsection{9}{物理の本は、もうこの巻の道具を使っている}

\label{bk:8-0:butsuri}物理シリーズ第1巻（力学）では、この巻の道具を、章の番号を書いて先に使っています。この巻を読むと、物理の本のその場所が「わかって」読めるようになります。

\begin{bkcenter}
\begin{bktable}{colspec={Q[c,m,wd=12mm] Q[l,m,wd=58mm] Q[l,m,wd=34mm] Q[c,m,wd=20mm]}, rowsep=1.2pt, colsep=3pt, cells={font=\small}, row{1}={font=\sffamily\bfseries\small, bg=bkTheme!10}}
物理 & そこで使っていること & この巻の道具 & 章 \\
力-4 & 力は、位置エネルギーが下がる向きに働く & 勾配 $\mathrm{grad}\,U$ & 8-4 \\
力-4 & 遠回りすると仕事が変わる力・変わらない力 & 線積分、回転 & 8-7・8-10 \\
力-6 & 回しにくさ（慣性モーメント） & 重積分 & 8-6 \\
力-7 & 関節を少し回すと、手はどれだけ動く？ & ヤコビ行列 $J$ & 8-3 \\
力-12 & 水の圧力・浮力 & 勾配・ガウスの定理 & 8-4・8-11 \\
力-13 & 流れこむ水と流れ出る水 & 面積分・発散 & 8-8・8-9 \\
力-15 & ねばりの力は「まわりの平均との差」 & ラプラシアン $\nabla^2$ & 8-12 \\
\end{bktable}
\bktabcaption{物理シリーズ第1巻で、この巻の道具を使っている所（一部）}
\end{bkcenter}

逆に、この巻の例題や「物理・工学ではこう使う」には、物理の本と同じ場面・同じ数を使っています（ロボットの腕の長さ、仕事の道A・道B など）。どちらから読んでも、同じ場面に出会えるようにしてあります。

\bksubsection{10}{くり返し出てくる場面}

\label{bk:8-0:bamen}主役の山と6つの場のほかにも、この巻では、同じ場面を、同じ数のまま、いくつもの章で使います。前の章で計算した答えが、あとの章で「別の方法でも同じになる」ことを確かめられるようにするためです。

\begin{bkcenter}
\begin{bktable}{colspec={Q[l,m,wd=27mm] Q[l,m,wd=66mm] Q[c,m,wd=27mm]}, rowsep=1.2pt, cells={font=\small}, row{1}={font=\sffamily\bfseries\small, bg=bkTheme!10}}
場面 & どんな場面か & 出てくる章 \\
ロボットの腕 & 2つの関節の腕（第7巻 7-2・物理 力-7 と同じ長さ） & 8-3 \\
厚さがちがう板 & 場所で厚さが変わる板の重さ & 8-6 \\
池 & 水がわき出す口と、吸いこまれる口がある池 & 8-5・8-9・8-11 \\
管の流れ & 水道管の中の、まん中ほど速い流れ & 8-6・8-8 \\
かべ・金属の棒 & 6-3 の熱が伝わるかべ。熱が広がる棒 & 8-4・8-12・8-13 \\
点電荷・電流・地球 & 電気の力・磁石の力・重力の矢印 & 8-5・8-8・8-11・8-14 \\
\end{bktable}
\bktabcaption{この巻でくり返し出てくる場面}
\end{bkcenter}

%% ---------------- ⑥ 表計算と Python ----------------
\bksection{doguu}{表計算と Python}

\bksubsection{11}{2次元の表は、地図になる}

\label{bk:8-0:dougu}場所ごとに決まる量は、数がとてもたくさんあります。そこで、この巻では、表計算（2-18）と Python（2-19）をたくさん使います。表計算の表は、たて・横にマスが並んでいるので、\textbf{そのまま地図}として使えます。この巻では、A列に北への距離 $y$、1行めに東への距離 $x$ を入れ、北が上になるように並べます（くわしくは 8-1）。

\begin{bkcenter}
\kviiiafx{B2}{=100-0.01*B\$1\^{}2-0.02*\$A2\^{}2}\par\vspace{1mm}
\bksheet{7}{12}{%
1 & \SetCell{c} $y$＼$x$ & \SetCell{bg=kviiiaX} 0 & \SetCell{bg=kviiiaX} 10 & \SetCell{bg=kviiiaX} 20 & \SetCell{bg=kviiiaX} 30 & \SetCell{bg=kviiiaX} 40 & \SetCell{bg=kviiiaX} 50 \\
2 & \SetCell{bg=kviiiaY} 40 & 68 & 67 & 64 & 59 & 52 & 43 \\
3 & \SetCell{bg=kviiiaY} 30 & 82 & 81 & 78 & 73 & 66 & 57 \\
4 & \SetCell{bg=kviiiaY} 20 & 92 & 91 & 88 & \SetCell{bg=kviiiaP} 83 & 76 & 67 \\
5 & \SetCell{bg=kviiiaY} 10 & 98 & 97 & 94 & 89 & 82 & 73 \\
6 & \SetCell{bg=kviiiaY} 0 & 100 & 99 & 96 & 91 & 84 & 75 \\}
\bktabcaption{主役の山の高さ（m）を表計算で地図にしたもの。黄色のマス E4 が、点 $(30,\ 20)$ の高さ 83 m。}
\end{bkcenter}

Python では、numpy で同じ表を作り、matplotlib という道具で、等高線の図（\texttt{contour}）や矢印の図（\texttt{quiver}）をかきます。この巻に出てくるプログラムは短いものばかりで、図はどれも本の形にかき直して並べてあります。

\Needspace{14\baselineskip}
\begin{bkcode}[Python]
import numpy as np               \# 数の表の道具 numpy を np として使う
x = np.arange(0, 60, 10)         \# 東へ 0, 10, …, 50（60 は入らない）
y = np.arange(40, -10, -10)      \# 北へ 40, 30, …, 0（北を上に）
X, Y = np.meshgrid(x, y)         \# マスごとの x の表と y の表
H = 100 - 0.01*X**2 - 0.02*Y**2  \# マスごとの高さ（B2 の式と同じ）
print(H)
\end{bkcode}
\begin{bkcodeout}[実行結果]
[[ 68.  67.  64.  59.  52.  43.]
 [ 82.  81.  78.  73.  66.  57.]
 [ 92.  91.  88.  83.  76.  67.]
 [ 98.  97.  94.  89.  82.  73.]
 [100.  99.  96.  91.  84.  75.]]
\end{bkcodeout}

\noindent 実行結果の表は、表計算の表とまったく同じ並びです。\texttt{68.} の「.」は、小数（68.0）として計算したという印です。

\begin{column}{ベクトル解析の記号は、だれが作った？}
この巻で学ぶ数学を、大学では「\textbf{ベクトル解析}」とよびます。場所ごとの矢印を微分したり積分したりする道具です。

19世紀の中ごろ、電気と磁気の研究が進み、場所ごとの矢印を計算する方法が必要になりました。はじめは、ハミルトンが考えた「四元数」という数を使って計算していました。19世紀の終わりごろ、アメリカのギブズと、イギリスのヘビサイドが、それぞれ別に、四元数から使いやすい部分だけをとり出して、今のベクトル解析の形にまとめました。1901年に、ギブズの講義をもとにしたウィルソンの本が出て、今使われている記号やことばの多くが、このころに決まりました。

8-4 で出てくる逆三角の記号 $\nabla$ は「ナブラ」と読みます。1870年11月10日に、イギリスのロバートソン・スミスが、物理学者テイトへの手紙の中で、この形が古代のたて琴（ギリシャ語で nabla）に似ていることから、この名前を提案したと伝えられています。
\end{column}

%% ---------------- ⑦ 新しい記号 ----------------
\bksection{card}{この巻で出会う新しい記号}

\label{bk:8-0:kigo}第8巻では、次のような新しい記号に出会います。いまは意味がわからなくてもだいじょうぶです。どの記号も、使う前に、数の例と図でていねいに説明します（記号カードは巻末にも並べてあります）。

\begin{bkcenter}
\begin{bktable}{colspec={Q[c,m,wd=36mm] Q[l,m,wd=30mm] Q[l,m,wd=46mm] Q[c,m,wd=9mm]}, rowsep=0.6pt, colsep=3pt, cells={font=\small}, row{1}={font=\sffamily\bfseries\small}}
記号 & 読み方 & ひとことで言うと & 章 \\
$f(x,\ y)$ & エフ・エックス・ワイ & 2つの数を入れると決まる数 & 8-1 \\
$\dfrac{\partial f}{\partial x}$ & ラウンド・ディー・エフ・ラウンド・ディー・エックス & ほかを止めて、$x$ だけ動かした変わり方（偏微分） & 8-2 \\
$\dd f$、$J$ & ディー・エフ、ジェイ & 少し動いたときの変化（全微分）、ヤコビ行列 & 8-3 \\
$\mathrm{grad}\,f$、$\nabla$ & グラディエント、ナブラ & いちばん急な向きの矢印（勾配） & 8-4 \\
$\vec F(x,\ y,\ z)$ & エフ & 場所ごとの矢印（ベクトル場） & 8-5 \\
$\displaystyle\iint_D$、$\displaystyle\iiint_V$ & 重積分、3重積分 & 面・立体のタイルでたす & 8-6 \\
$\displaystyle\int_C\vec F\cdot\dd\vec r$、$\displaystyle\oint_C$ & 線積分 & 道にそってたす。◯は閉じた道 & 8-7 \\
$\vec n\,\dd S$、$\displaystyle\oiint_S$ & エヌ・ディー・エス & 向きのついた面のタイル。◯は閉じた面 & 8-8 \\
$\mathrm{div}\,\vec F$ & ダイバージェンス & わき出しの量（発散） & 8-9 \\
$\mathrm{rot}\,\vec F$ & ローテーション & うずの強さ（回転） & 8-10 \\
$\nabla^2 f$ & ナブラ2乗 & まわりの平均との差（ラプラシアン） & 8-12 \\
$\dfrac{\mathrm D T}{\mathrm D t}$ & ディー・ティー・ディー・ティー & 流れといっしょに動いて見た変化 & 8-13 \\
\end{bktable}
\bktabcaption{この巻で出会う記号（くわしくはそれぞれの章で）}
\end{bkcenter}

この巻では、ベクトルは第4巻・第7巻と同じく、矢印つきの文字（$\vec F$）で書きます。大学の本では、太字（$\symbf F$）で書くことが多いので、それぞれの章の「専門書をのぞいてみよう」と、8-14 で見くらべます。

%% ---------------- ⑧ 路線 ----------------
\bksection{kviiiaroute}{この巻が通る路線と終点}

\label{bk:8-0:rosen}第8巻を通る路線は3本です。3本とも、この巻に\textbf{終点}があります。ここまで小学校から乗ってきた路線が、この巻でゴールに着き、物理の世界（流体・熱・電磁気・ロボット）へつながります。

\begin{bkfigure}
\begin{tikzpicture}[x=1mm, y=1mm]
  \kviiiaroute{28}{1}{6-1}{8-1〜8-3・{\bfseries 8-4}・8-12・8-13}{終点}
  \kviiiaroute{14}{2}{6-3}{8-6〜8-8・{\bfseries 8-11}}{終点}
  \kviiiaroute{0}{6}{7-7}{{\bfseries 8-5}・8-9・8-10・8-14}{終点}
  \node[font=\sffamily\bfseries\small, text=bkGray] at (24,38) {前の駅};
  \node[font=\sffamily\bfseries\small, text=bkDeep] at (74,38) {この巻（第8巻）};
  \node[font=\sffamily\bfseries\small, text=bkGray] at (126,38) {};
\end{tikzpicture}
\bkfigcaption{第8巻を通る3本の路線。四角の中がこの巻の駅（太字が終点の駅）。右はしの四角が終点の印。終点は、その路線の大きな目的地。終点のあとの駅（8-12 など）は、終点で作った道具を使う駅。}
\end{bkfigure}

\begin{bkcenter}
\begin{bktable}{colspec={Q[c,m,wd=7mm] Q[l,m,wd=20mm] Q[l,m,wd=94mm]}, rowsep=0.8pt, cells={font=\small}, row{1}={font=\sffamily\bfseries\small}}
 & 路線 & 大きな駅のつながり \\
\bkrouteicon{1} & 変化 & 傾き（3-5）→ 微分（5-1）→ 微分方程式（6-1）→ \textbf{勾配}（8-4、終点） \\
\bkrouteicon{2} & たし合わせ & 細かく分けてたす（4-5）→ 定積分（5-10）→ 積分で解く（6-3）→ \textbf{ガウスの定理}（8-11、終点） \\
\bkrouteicon{6} & 矢印 & 座標（3-2）→ ベクトル（4-10）→ 3次元の回転（7-7）→ \textbf{ベクトル場}（8-5、終点） \\
\end{bktable}
\bktabcaption{3本の路線の、大きな駅}
\end{bkcenter}

どの路線も、第1巻・第2巻の章から始まっています（変化は 1-5「ひき算」、たし合わせは 1-21「広さ」、矢印は 2-2「数直線」）。

%% ---------------- この巻の読み方 ----------------
\Needspace{27\baselineskip}\bksection{hon}{この巻の読み方}

\label{bk:8-0:yomikata}
\begin{bkfigure}
\begin{tikzpicture}[x=1mm, y=1mm]
  \kviiiatalk{0}{sora}{ソラ}{第8巻は、大学2年までの内容なんだね。$\partial$ とか $\nabla$ とか、見たことのない記号がいっぱい。}
  \kviiiatalk{-22}{mina}{ミナ}{でも、地図の等高線や、天気図の風の矢印なら、見たことがあるよ。それを数で読むんだって。}
  \kviiiatalk{-46}{delta}{デルタ}{そのとおり。1回読んで分からなくても、だいじょうぶ。図をかいて、数の表で手を動かせば、かならずわかるよ。}
\end{tikzpicture}
\end{bkfigure}

\begin{bkfigure}
\begin{tikzpicture}[x=1mm, y=1mm]
  \kviiiacard{(0,0)}{shiki}{bkNavy}{①\enspace 成分で書いてから記号に}
    {記号 $\mathrm{grad}$・$\mathrm{div}$・$\mathrm{rot}$ だけで考えず、まず成分で1つずつ書きます。\par\vspace{2pt}
     \kviiiaitem{$x$ の成分、$y$ の成分を1行ずつ。}\kviiiaitem{記号の式のあとには、主役の山や主役の場の数で1回確かめる。}}
  \kviiiacard{(68,0)}{me}{bkGreen}{②\enspace 図と数の表の2通りで}
    {初めての考えは、図と数の表の2通りで見せます。\par\vspace{2pt}
     \kviiiaitem{等高線・矢印・タイルの図を見る。}\kviiiaitem{同じことを、表計算のマスの数で確かめる。}}
  \kviiiacard{(0,-57)}{te}{bkOrange}{③\enspace 名前と、実際にすること}
    {積分の記号は「操作の名前」です。\par\vspace{2pt}
     \kviiiaitem{名前の段・直す段・楽をする段の3つで考える。}\kviiiaitem{場合分けの表で「この場合はこう直す」を見つける。}}
  \kviiiacard{(68,-57)}{goal}{bkDeep}{④\enspace ★5の章は3段}
    {大学の内容の章は、3つの段に分けてあります。\par\vspace{2pt}
     \kviiiaitem{「目でわかる」「手でわかる」「式でわかる」。}\kviiiaitem{「ここまでで十分」の旗まで読めば、その章のゴール。}}
\end{tikzpicture}
\end{bkfigure}
